\documentclass[10pt,a4paper]{amsart}
\usepackage[margin=19mm]{geometry}
\usepackage{amsmath,amssymb,mathtools}
\usepackage[scr=rsfs,cal=euler]{mathalfa}
\usepackage{xfrac}
\usepackage[unicode,hidelinks,bookmarks=false]{hyperref}
\newcommand{\ie}{i.e.\ }
\newcommand{\cf}{cf.\ }
\newcommand{\resp}{respectively\ }
\newcommand{\Subsection}{\subsection}
\providecommand{\nobreakdash}{\nobreak\hbox{-}}
\setlength{\emergencystretch}{2em}
\multlinegap0pt
\usepackage{amssymb}
\usepackage{bm}
\usepackage{mathtools}
\newtheorem{theorem}{Theorem}
\newtheorem{corollary}{Corollary}
\newtheorem{lemma}{Lemma}
\newtheorem{proposition}{Proposition}
\usepackage{cleveref}
\crefname{theorem}{Theorem}{Theorems}
\crefname{corollary}{Corollary}{Corollarys}
\crefname{lemma}{Lemma}{Lemmas}
\crefname{proposition}{Proposition}{Propositions}
\newcommand{\LC}{\mathcal{L}}
\newcommand{\LO}{\mathrm{LO}} % Lieb--Oxford
\newcommand{\Lloc}{L_{\mathrm{loc}}} % Locally integrable
\newcommand{\NUEG}{\mathrm{NUEG}} % Non-uniform electron gas
\newcommand{\RR}{\mathbb{R}}
\newcommand{\SO}{\mathrm{SO}}
\newcommand{\TF}{\mathrm{TF}} % Thomas--Fermi
\newcommand{\UEG}{\mathrm{UEG}} % Uniform electron gas
\def\XXint#1#2#3{{\setbox0=\hbox{$#1{#2#3}{\int}$ }
\vcenter{\hbox{$#2#3$ }}\kern-.6\wd0}}
\def\Xint#1{\mathchoice
{\XXint\displaystyle\textstyle{#1}}%
{\XXint\textstyle\scriptstyle{#1}}%
{\XXint\scriptstyle\scriptscriptstyle{#1}}%
{\XXint\scriptscriptstyle\scriptscriptstyle{#1}}%
\!\int}
\newcommand{\ZZ}{\mathbb{Z}}
\def\dashint{\Xint-}
\newcommand{\dd}{\mathrm{d}} % Differential
\renewcommand{\epsilon}{\varepsilon}
\renewcommand{\ge}{\geqslant}
\newcommand{\grad}{\bm{\nabla}} % Gradient
\newcommand{\iden}{\mathbb{1}} % Indicator, identity
\renewcommand{\le}{\leqslant}
\newcommand{\ol}[1]{\overline{#1}}
\newcommand{\tetra}{\mathbb{\Delta}} % Tetrahedron
\newcommand{\ul}[1]{\underline{#1}}
\newcommand{\wt}[1]{\widetilde{#1}}
\title{The non-uniform electron gas}
\author{Mihaly A. Csirik, Andre Laestadius}
\date{Source: arXiv:2603.21200v1 (CC BY 4.0)}
\begin{document}
\maketitle

\noindent\emph{Source and licence.} The non-uniform electron gas. Version arXiv:2603.21200v1 (CC BY 4.0), distributed by arXiv under the Creative Commons Attribution 4.0 International licence, http://creativecommons.org/licenses/by/4.0/, as recorded in the arXiv licence metadata for this exact version and shown on the abstract page https://arxiv.org/abs/2603.21200v1. Full licence notice: \texttt{licenses/arxiv-2603.21200-CC-BY-4.0.txt}.\par\smallskip

\noindent\textit{Editorial scope.} Counted unit: the article's theorem tetrathermo (source0.tex lines 1405--1429). The article's complete original proof of it is delivered with it (source0.tex lines 1526--1640, one \texttt{proof} environment of 115 lines, which the article places away from the statement). No mathematics is written, repaired or re-derived: the statement and the proof are the article's own lines, in the article's own notation, and no retained line is substituted or re-flowed.  Retained uncounted as the source-local context the proof needs: lines 262--268; lines 1089--1095; lines 1128--1135; lines 1175--1189; lines 1287--1301; lines 1436--1442; lines 1518--1520.  Nothing was omitted from any retained span, so the package carries no omission record.  No retained line contains a citation, so the wrapper carries no bibliography and no bibliography entry.  No retained line contains a figure environment, an image-inclusion command or drawing code, so no image is delivered and the package declares no asset dependency.  Editorial formatting only, in addition: an AMS \texttt{amsart} preamble that restates the article's own declarations and macros used by the retained lines, namely the environments \texttt{theorem}, \texttt{corollary}, \texttt{lemma}, \texttt{proposition} on separate counters, together with the standard packages those lines need.  The retained environments print in this document as Theorem 1, Corollary 1, Lemma 1, Proposition 1 in order of appearance, whereas the article prints the same items with the numbering of its own full document.  The complete native source of the article as distributed in the arXiv e-print for this exact version is bundled unchanged as \texttt{sources/196925/source0.tex}.
\begin{theorem}\label{permeanthm}
For any $\LC$-periodic complex-valued function $u\in\Lloc^1(\RR^d)$, the mean value of $u$ over $\RR^d$ exists and equals its mean value over a unit cell $\Lambda\subset\RR^d$ of $\LC$, i.e.
\begin{equation}\label{permean}
\dashint_{\RR^d} u:=\lim_{L\to\infty} \dashint_{C_L+a} u =\dashint_{\Lambda} u,
\end{equation}
independently and uniformly in $a\in\RR^d$.
\end{theorem}

\begin{corollary}\label{llsebound}
Suppose that $\rho\in L^1(\RR^3;\RR_+)$ and $\sqrt{\rho}\in H^1(\RR^3)$. Then for every $0<\epsilon\le 1/15$ the bound
$$
E(\rho)\le (1+\epsilon)c_\TF\hbar^2 \int_{\RR^3} \rho^{5/3} + \frac{19}{15}\frac{\hbar^2}{\epsilon} \int_{\RR^3} |\grad\sqrt{\rho}|^2
$$
holds true. 
\end{corollary}

\begin{lemma}[Bound on kinetic energy of a smeared set]\label{indgradlemma}
Suppose that $\Omega\subset\RR^3$ is an open set with a boundary having finite 2-dimensional Lebesgue measure: $|\partial\Omega|<\infty$. 
Then the following bound holds true
$$
\int_{\RR^3} \left|\grad \sqrt{\iden_\Omega*\eta_\delta}\right|^2\le \frac{C_\eta|\partial\Omega|}{\delta},
$$
where the constant $C_\eta>0$ depends only on the regularization function $\eta$.
\end{lemma}

\begin{proposition}\label{basicbdprop}
The following a priori bounds hold true.
\begin{itemize}
\item[(i)] For every $0<\epsilon\le 1/15$ and every inhomogeneity $\zeta$, there holds
\begin{multline*}
f_{\Omega,\eta_\delta,\zeta}(a)\le (1+\epsilon) c_{\TF}\hbar^2  \frac{1}{|\Omega|} \int_{\RR^3} (\iden_{\Omega}*\eta_\delta)\zeta^{5/3}(\cdot-a) \\
+ \frac{38}{15}\frac{\hbar^2}{\epsilon} \frac{1}{|\Omega|} \int_{\RR^3} (\iden_{\Omega}*\eta_\delta)|\grad \sqrt{\zeta}|^2(\cdot-a)+ \frac{38}{15}\frac{\hbar^2}{\epsilon} \frac{1}{|\Omega|} \int_{\RR^3} |\grad \sqrt{\iden_{\Omega}*\eta_\delta}|^2\zeta(\cdot-a) .
\end{multline*}
\item[(ii)] For every $0<\epsilon\le 3/5$ and every inhomogeneity $\zeta$, there holds 
\begin{multline*}
f_{\Omega,\eta_\delta,\zeta}(a)\ge (1-\epsilon) c_{\TF}\hbar^2  \frac{1}{|\Omega|} \int_{\RR^3} (\iden_{\Omega}*\eta_\delta)\zeta^{5/3}(\cdot-a) -c_\LO \frac{1}{|\Omega|}\int_{\RR^3} (\iden_{\Omega}*\eta_\delta)\zeta^{4/3}(\cdot-a) \\
- \frac{20}{27}\frac{\hbar^2}{\epsilon} \frac{1}{|\Omega|} \int_{\RR^3} (\iden_{\Omega}*\eta_\delta)|\grad \sqrt{\zeta}|^2(\cdot-a) + \frac{20}{27}\frac{\hbar^2}{\epsilon} \frac{1}{|\Omega|} \int_{\RR^3} |\grad \sqrt{\iden_{\Omega}*\eta_\delta}|^2\zeta(\cdot-a) .
\end{multline*}
\end{itemize}
\end{proposition}

\begin{theorem}[Improved decoupling upper bound]\label{impdecup}
Let $0<\delta\le \ell/2$ and $0<\alpha<1/2$. Then for any $\rho\in L^1(\RR^3,\RR_+)$ with $\grad\sqrt{\rho}\in L^2(\RR^3)$ the bounds
\begin{multline*}
E(\rho)\le  \dashint_{1-\alpha}^{1+\alpha}\frac{\dd t}{t^4} \int_{\SO(3)}\dd R \dashint_{C_{t\ell}}\dd\tau \sum_{z\in\ZZ^3}  \sum_{j=1}^{24} 
E\bigl(\iden_{t\ell T_j(1-\delta/\ell)\tetra} * \eta_\delta(R\cdot - t\ell z - \tau)\rho\bigr)\\
+\frac{C\hbar^2\delta}{\ell} \int_{\RR^3}\rho^{5/3} +  \frac{C\hbar^2}{\ell\delta} \int_{\RR^3} \rho + \frac{C\hbar^2\delta}{\ell}\int_{\RR^3} |\grad\sqrt{\rho}|^2
+C\delta^2\log(\alpha^{-1}) \int_{\RR^3} \rho^2
\end{multline*}
and 
\begin{multline*}
T(\rho)\le \dashint_{C_{\ell}}\dd\tau \sum_{z\in\ZZ^3}  \sum_{j=1}^{24} T\bigl(\iden_{\ell T_j(1-\delta/\ell)\tetra} * \eta_\delta(\cdot - \ell z - \tau)\rho\bigr)\\
+\frac{C\hbar^2\delta}{\ell} \int_{\RR^3}\rho^{5/3} +  \frac{C\hbar^2}{\ell\delta} \int_{\RR^3} \rho + \frac{C\hbar^2\delta}{\ell}\int_{\RR^3} |\grad\sqrt{\rho}|^2
\end{multline*}
hold true for some universal constant $C>0$.
\end{theorem}

\begin{equation}\label{tetracomp}
\begin{multlined}
e_{\ell'\tetra,\eta_{\delta'}}(\zeta)\ge \Bigg( 1-C\sigma \frac{\delta}{\ell} -C\sigma\frac{\ell+\delta+\delta'}{\ell'}\Bigg) e_{\ell\tetra,\eta_{\delta}}(\zeta)
- C \frac{\ell+\delta+\delta'}{\ell'} \dashint_{\RR^3} \zeta^{4/3}\\
-\frac{C}{\ell}\Bigg( (1+\hbar^2\delta^{-1}) \dashint_{\RR^3} \zeta + \delta^3\dashint_{\RR^3} \zeta^2 \Bigg)
\end{multlined}
\end{equation}

\begin{equation}\label{volbdbound}
\frac{M_\ell^\partial}{|\ell'\tetra|}\le C \frac{\ell+\delta+\delta'}{\ell'}.
\end{equation}

\begin{theorem}[Convergence rate for tetrahedra]\label{tetrathermo}
Fix an inhomogeneity $\zeta$ and a regularization function $\eta$. Let $\ul{e}_\NUEG(\zeta)$ and $\ol{e}_\NUEG(\zeta)$ denote the liminf and the limsup of $e_{\ell\tetra,\eta_\delta}(\zeta)$
as $\ell\delta\to\infty$ and $\delta/\ell\to 0$. 
\begin{itemize}
\item[(i)] For $\ell/\delta$ sufficiently large, there holds
\begin{align*}
&e_{\ell\tetra,\eta_{\delta}}(\zeta)\le \ul{e}_\NUEG(\zeta) +\frac{C}{\ell}\Bigg( (1+\hbar^2\delta^{-1}) \dashint_{\RR^3} \zeta + \delta^3\dashint_{\RR^3} \zeta^2 + \hbar^2\delta \dashint_{\RR^3} \zeta^{5/3} + \hbar^2\delta \dashint_{\RR^3} |\grad\sqrt{\zeta}|^2\Bigg)
\end{align*}
\item[(ii)] For $\ell/\delta$ sufficiently large and $0<\alpha<1/2$, 
\begin{multline*}
\ol{e}_{\NUEG}(\zeta)\le \dashint_{1-\alpha}^{1+\alpha}\frac{\dd t}{t^4}\, e_{t\ell\tetra,\eta_{t\delta}}(\zeta)\\
+\frac{C\hbar^2\delta}{\ell} \dashint_{\RR^3}\zeta^{5/3} +  \frac{C\hbar^2}{\ell\delta} \dashint_{\RR^3} \zeta + \frac{C\hbar^2\delta}{\ell}\dashint_{\RR^3} |\grad\sqrt{\zeta}|^2+ C\delta^2\log(\alpha^{-1}) \dashint_{\RR^3}\zeta^2
\end{multline*}
\item[(iii)] For $\ell$  sufficiently large, 
\begin{align*}
&e_{\ell\tetra,\eta_{\delta}}(\zeta)\ge \ol{e}_{\NUEG}(\zeta) -\frac{C\delta}{\ell}\dashint_{\RR^3} \left(\zeta +\zeta^2\right) -\frac{C}{\ell^{2/5}}\dashint_{\RR^3} \left(\zeta +\zeta^2\right) - \frac{C}{\ell^{4/5}}\dashint_{\RR^3} |\grad\sqrt{\zeta}|^2.
\end{align*}
\item[(iv)] The thermodynamic limit
$$
\lim_{\substack{\delta/\ell\to 0\\\ell\delta\to\infty}} e_{\ell\tetra,\eta_\delta}(\zeta)=e_\NUEG(\zeta)
$$
exists and is independent of the regularization function $\eta$.
\end{itemize}
In all the above bounds, the generic constant $C>0$ depends only on the regularization function $\eta$ and is independent of $\ell$, $\delta$, $\zeta$ and $\alpha$. 
\end{theorem}

\begin{proof}[Proof of \cref{tetrathermo}]
\noindent\textbf{Proof of (i).}
We split the l.h.s. of \cref{tetracomp} according to 
\begin{align*}
e_{\ell'\tetra,\eta_{\delta'}}(\zeta)&=(1-C\sigma\delta/\ell) e_{\ell'\tetra,\eta_{\delta'}}(\zeta) +  \frac{C\sigma\delta}{\ell} e_{\ell'\tetra,\eta_{\delta'}}(\zeta)\\
&\le (1-C\sigma\delta/\ell) e_{\ell'\tetra,\eta_{\delta'}}(\zeta) + \frac{C\hbar^2\sigma\delta}{\ell} \Bigg[ \dashint_{\RR^3} \zeta^{5/3} + \dashint_{\RR^3} |\grad\sqrt{\zeta}|^2 +  \frac{1}{\ell'\delta'}\dashint_{\RR^3} \zeta \Bigg]
\end{align*}
where we used \cref{basicbdprop} (i) with $\epsilon=1$. Rearranging our bound, we find
\begin{multline*}
\Bigg( 1-C\sigma \frac{\delta}{\ell} -C\sigma\frac{\ell+\delta+\delta'}{\ell'}\Bigg) e_{\ell\tetra,\eta_{\delta}}(\zeta)\\
\le (1-C\sigma\delta/\ell) e_{\ell'\tetra,\eta_{\delta'}}(\zeta)
 +\frac{C}{\ell}\Bigg( (1+\hbar^2\delta^{-1}) \dashint_{\RR^3} \zeta + \delta^3\dashint_{\RR^3} \zeta^2 + \hbar^2\delta \dashint_{\RR^3} \zeta^{5/3} + \hbar^2\delta \dashint_{\RR^3} |\grad\sqrt{\zeta}|^2\Bigg) \\
 + C \frac{\ell+\delta+\delta'}{\ell'} \dashint_{\RR^3} \zeta^{4/3} + \frac{C\hbar^2\sigma\delta}{\ell}  \frac{1}{\ell'\delta'}\dashint_{\RR^3} \zeta
\end{multline*}
Now taking the liminf as $\ell'\to\infty$, and dividing by $(1-C\sigma\delta/\ell)$, we arrive at
\begin{align*}
&e_{\ell\tetra,\eta_{\delta}}(\zeta)\le \ul{e}_\UEG(\zeta) +\frac{C}{\ell}\Bigg( (1+\hbar^2\delta^{-1}) \dashint_{\RR^3} \zeta + \delta^3\dashint_{\RR^3} \zeta^2 + \hbar^2\delta \dashint_{\RR^3} \zeta^{5/3} + \hbar^2\delta \dashint_{\RR^3} |\grad\sqrt{\zeta}|^2\Bigg),
\end{align*}
which is what we wanted to show.

%%%%%%%%%%%%%%%%%%%%%%%%%%%%%%%%%%%%%%%%%%%%%%%%%%%%%%%%%%%%%%%%%%%%%%%%%%%%%%%%%%%%%%%%
\vspace{.3em}
\noindent\textbf{Proof of (ii).} The proof of the lower bound is similar. Using \cref{impdecup}, we have for $0<\alpha<\frac{1}{2}$, 
\begin{equation}\label{etetraup}
\begin{aligned}
&e_{\ell'\tetra,\eta_{\delta'}}(\zeta)\le\frac{1}{|\ell'\tetra|}\int_{\SO(3)}\dd Q \dashint_{\RR^3}\dd a \dashint_{1-\alpha}^{1+\alpha}\frac{\dd t}{t^4} \int_{\SO(3)}\dd R 
\dashint_{C_{t(\ell+\delta)}}\dd\tau\\
	&\times \sum_{z\in\ZZ^3} \sum_{j=1}^{24} E\bigl((\iden_{t(\ell+\delta) T_j \beta\tetra} * \eta_{t \delta})(R\cdot - t(\ell+\delta) z - \tau)(\iden_{\ell'\tetra}*\eta_{\delta'})\zeta(Q(\cdot - a))\bigr)\\
	&+\Biggl[\frac{C\hbar^2\delta}{\ell} \dashint_{\RR^3}\zeta^{5/3} +  \frac{C\hbar^2}{\ell\delta} \dashint_{\RR^3} \zeta + \frac{C\hbar^2\delta}{\ell}\dashint_{\RR^3} |\grad\sqrt{\zeta}|^2 + C\delta^2\log(\alpha^{-1}) \dashint_{\RR^3}\zeta^2\Biggr]\\
    &=\mathrm{(I)} + \mathrm{(II)}
\end{aligned}
\end{equation}
where we used the notation $\beta=\ell/(\ell+\delta)$, where $\frac{2}{3}\le\beta<1$. Notice that the scales are chosen so that the small tetrahedra are of size $t\ell$ and smeared at scale $t\delta$.

Again, we split (I) according to $J_0$ and $J\setminus J_0$, which we define analogously to above.
First, we consider the summation over $J_0$ in (I), which we denote by (Ia).
Using $t(\ell+\delta) T_j(1-\delta/(\ell+\delta))\tetra=t(\ell+\delta)\beta R_j\tetra + t(\ell+\delta) z_j$, we may write
\begin{multline*}
(\iden_{t(\ell+\delta) T_j\beta\tetra} * \eta_{t\delta})(Rx - t(\ell+\delta) z - \tau)\\
=(\iden_{t\ell\tetra} * \eta_{t\delta})(R_j^\top Rx - R_j^\top(t(\ell+\delta) z + \tau + t(\ell+\delta) z_j))
\end{multline*}
We have
\begin{align*}
&\mathrm{(Ia)}=\frac{1}{|\ell'\tetra|} \dashint_{1-\alpha}^{1+\alpha}\frac{\dd t}{t^4}  \int_{\SO(3)}\dd Q \dashint_{\RR^3}\dd a \int_{\SO(3)}\dd R \dashint_{C_{t(\ell+\delta)}}\dd\tau\sum_{(z,j)\in J_0}\\
	&\times \dashint_{\RR^3}\int_{\SO(3)} E\bigl((\iden_{t\ell\tetra} * \eta_{t\delta})(R_j^\top R\cdot - R_j^\top(t(\ell+\delta) z + \tau + t(\ell+\delta) z_j))\zeta(Q(\cdot - a))\bigr)\\
&=\frac{1}{|\ell'\tetra|} \dashint_{1-\alpha}^{1+\alpha}\frac{\dd t}{t^4} |J_0|
	\int_{\SO(3)}\dd Q \dashint_{\RR^3}\dd a\,  E\bigl((\iden_{t\ell\tetra} * \eta_{t\delta})\zeta( Q(\cdot- a))\bigr)\\
	&=\dashint_{1-\alpha}^{1+\alpha}\frac{\dd t}{t^4} \frac{M_{t\ell}^0}{|\ell'\tetra|} e_{t\ell\tetra,\eta_{t\delta}}(\zeta)
\end{align*}
In the second equality, we used the isometry-invariance of the energy and \cref{permeanthm}.
Here, the volume ratio may be bounded similarly as before
$$
\frac{M_{t\ell}^0}{|\ell'\tetra|}\le 1+C\sigma\frac{t\ell+t\delta+\delta'}{\ell'}\le 1+C\sigma\frac{\ell+\delta+\delta'}{\ell'}.
$$
Next, we estimate the contributions near the boundary using \cref{llsebound},
\begin{multline*}
\mathrm{(Ib)}=\frac{1}{|\ell'\tetra|}\int_{\SO(3)}\dd Q \dashint_{\RR^3}\dd a \dashint_{1-\alpha}^{1+\alpha}\frac{\dd t}{t^4} \int_{\SO(3)}\dd R \dashint_{C_{t(\ell+\delta)}}\dd\tau
\sum_{(z,j)\in J\setminus J_0} \times\\
	\times E\bigl((\iden_{t(\ell+\delta) T_j \beta\tetra} * \eta_{t\delta})(R\cdot - t(\ell+\delta) z - \tau)(\iden_{\ell'\tetra}*\eta_{\delta'})\zeta(Q(\cdot - a))\bigr)\\
\le \hbar^2\frac{1}{|\ell'\tetra|}\int_{\SO(3)}\dd Q \dashint_{\RR^3}\dd a \dashint_{1-\alpha}^{1+\alpha}\frac{\dd t}{t^4} \int_{\SO(3)}\dd R \dashint_{C_{t(\ell+\delta)}}\dd\tau
\sum_{(z,j)\in J\setminus J_0} \times\\
	\times\Bigg[ C\int_{\RR^3}  (\iden_{t(\ell+\delta) T_j \beta\tetra} * \eta_{t\delta})(R\cdot - t(\ell+\delta) z - \tau)(\iden_{\ell'\tetra}*\eta_{\delta'})\zeta(Q(\cdot - a))^{5/3}\\
    + C \int_{\RR^3} \Bigl|\grad\sqrt{(\iden_{t(\ell+\delta) T_j \beta\tetra} * \eta_{t\delta})(R\cdot - t(\ell+\delta) z - \tau)(\iden_{\ell'\tetra}*\eta_{\delta'})\zeta(Q(\cdot - a))}\Bigr|^2 \Bigg]
\end{multline*}
Here, the gradient term may be bounded using the Cauchy--Schwarz inequality as
\begin{align*}
C\int_{\RR^3} \Bigl|\grad\sqrt{(\iden_{t(\ell+\delta) T_j \beta\tetra} * \eta_{t\delta})(R\cdot - t(\ell+\delta) z - \tau)}\Bigr|^2(\iden_{\ell'\tetra}*\eta_{\delta'})\zeta(Q(\cdot - a))\\
+C\int_{\RR^3} \Bigl|\grad\sqrt{\iden_{\ell'\tetra}*\eta_{\delta'}}\Bigr|^2
(\iden_{t(\ell+\delta) T_j \beta\tetra} * \eta_{t\delta})(R\cdot - t(\ell+\delta) z - \tau)\zeta(Q(\cdot - a))\\
+C\int_{\RR^3} \bigl|\grad\sqrt{\zeta(Q(\cdot - a))}\bigr|^2(\iden_{t(\ell+\delta) T_j \beta\tetra} * \eta_{t\delta})(R\cdot - t(\ell+\delta) z - \tau)(\iden_{\ell'\tetra}*\eta_{\delta'})
\end{align*}
After computing the averages with respect to $a$ and $Q$, we obtain by \cref{indgradlemma},
\begin{equation}\label{ebduest}
\mathrm{(Ib)}\le C\hbar^2 \frac{M_{t\ell}^\partial}{|\ell'\tetra|}\Bigg[ \dashint_{\RR^3} \zeta^{5/3} + \frac{1}{\ell\delta}\dashint_{\RR^3}\zeta + \dashint_{\RR^3} |\grad\sqrt{\zeta}|^2\Bigg] + \frac{C\hbar^2}{\ell'\delta'}\dashint_{\RR^3}\zeta.
\end{equation}
Combining this with the bound \cref{volbdbound}, we arrive at an estimate on (Ib) that vanishes in the thermodynamic limit $\ell'\delta'\to\infty$, $\delta'/\ell'\to 0$. 

Collecting our estimates and taking the limsup as $\ell'\delta'\to\infty$, $\delta'/\ell'\to 0$, we obtain
\begin{multline*}
\ol{e}_{\NUEG}(\zeta)\le \dashint_{1-\alpha}^{1+\alpha}\frac{\dd t}{t^4} e_{t\ell\tetra,\eta_{t\delta}}(\zeta)\\
+\frac{C\hbar^2\delta}{\ell} \dashint_{\RR^3}\zeta^{5/3} +  \frac{C\hbar^2}{\ell\delta} \dashint_{\RR^3} \zeta + \frac{C\hbar^2\delta}{\ell}\dashint_{\RR^3} |\grad\sqrt{\zeta}|^2+ C\delta^2\log(\alpha^{-1}) \dashint_{\RR^3}\zeta^2
\end{multline*}
as stated.
%%%%%%%%%%%%%%%%%%%%%%%%%%%%%%%%%%%%%%%%%%%%%%%%%%%%%%%%%%%%%%%%%%%%%%%%%%%%%%%%%%%%%%%%

\vspace{.3em}
\noindent\textbf{Proof of (iii).} In order to get rid of the averaging in (ii), we use \cref{tetracomp} with $\ell\to t\ell$, $\delta\to t\delta$ and average over $t\in(1/2,3/2)$. We obtain
\begin{multline}\label{monoavg}
e_{\ell'\tetra,\eta_{\delta'}}(\zeta)\ge \Bigg( 1-C\sigma \frac{\delta}{\ell} -C\sigma\frac{\ell+\delta+\delta'}{\ell'}\Bigg) \dashint_{1-\alpha}^{1+\alpha}\frac{\dd t}{t^4} e_{t\ell\tetra,\eta_{t\delta}}(\zeta)\\
- C \frac{\ell+\delta+\delta'}{\ell'} \dashint_{\RR^3} \zeta^{4/3}-\frac{C}{\ell}\Bigg( (1+\hbar^2\delta^{-1}) \dashint_{\RR^3} \zeta + \delta^3\dashint_{\RR^3} \zeta^2 \Bigg).
\end{multline}
Using (ii), we may bound the $t$-average from below by 
$$
\ol{e}_{\NUEG}(\zeta) -\frac{C\hbar^2\delta}{\ell} \dashint_{\RR^3}\zeta^{5/3} -  \frac{C\hbar^2}{\ell\delta} \dashint_{\RR^3} \zeta - \frac{C\hbar^2\delta}{\ell}\dashint_{\RR^3} |\grad\sqrt{\zeta}|^2 -C\delta^2\dashint_{\RR^3}\zeta^2,
$$
to get after discarding the positive terms and rearranging,
\begin{align*}
&e_{\ell'\tetra,\eta_{\delta'}}(\zeta)\ge \ol{e}_{\NUEG}(\zeta) - C\frac{\ell+\delta+\delta'}{\ell'} \dashint_{\RR^3} \zeta^{4/3} \\
&-\frac{C}{\ell}\Biggl( (1+\hbar^2\delta^{-1}) \dashint_{\RR^3} \zeta + \delta^3\dashint_{\RR^3} \zeta^2 + \hbar^2\delta \dashint_{\RR^3} \zeta^{5/3} + \hbar^2\delta \dashint_{\RR^3} |\grad\sqrt{\zeta}|^2\Biggr) - C\delta^2 \dashint_{\RR^3}\zeta^2  \\
&\ge \ol{e}_{\NUEG}(\zeta) - C\frac{\ell+\delta+\delta'}{\ell'} \dashint_{\RR^3} \left(\zeta +\zeta^2\right) \\
&-C\frac{1+\delta^{-1}+\delta}{\ell} \dashint_{\RR^3} \zeta - C\left(\frac{\delta+\delta^3}{\ell}+\delta^2 \right)\dashint_{\RR^3} \zeta^2  -\frac{C\delta}{\ell} \dashint_{\RR^3} |\grad\sqrt{\zeta}|^2
\end{align*}
where in the last step we made use of $\zeta^\alpha\le \zeta+\zeta^2$ for $1\le\alpha\le 2$
Here, we choose $\delta=\ell^{-1/3}$ and then $\ell=(\ell')^{3/5}$ to obtain
\begin{align*}
&e_{\ell'\tetra,\eta_{\delta'}}(\zeta)\ge \ol{e}_{\NUEG}(\zeta) -\frac{C\delta'}{\ell'}\dashint_{\RR^3} \left(\zeta +\zeta^2\right) -\frac{C}{(\ell')^{2/5}}\dashint_{\RR^3} \left(\zeta +\zeta^2\right) - \frac{C}{(\ell')^{4/5}}\dashint_{\RR^3} |\grad\sqrt{\zeta}|^2
\end{align*}
for sufficiently large $\ell'$.

\vspace{.3em}
\noindent\textbf{Proof of (iv).} It only remains to explain why the limit is independent of the
regularization function $\eta$. Suppose that $\wt{\eta}$ is another regularization function. 
One can derive estimates similar to (i)--(iii) with finite-volume quantities involving $\wt{\eta}$ instead.
 This implies that the thermodynamic limit of $e_{\ell\tetra,\wt{\eta}_\delta}(\zeta)$ is the same as that of $e_{\ell\tetra,\eta_\delta}(\zeta)$.
\end{proof}

\end{document}
